\subsection*{Optical word meanings and literary source overlap}
We trained an English word-image encoder on 29,447 alphabetic terms using centred, whitened and unit-normalized 50-dimensional GloVe vectors.\footnote{\url{https://nlp.stanford.edu/projects/glove/}} Each input contained up to 24 occupied character positions within a word, including punctuation; visible blank cells delimited words. We then froze this encoder and fitted a Chinese predicate projection and an English question projection. A three-position convolution and learned attention pooled the optical word vectors. Given projected predicate $u$ and pooled question $v$, the semantic contribution to the score was $0.65\langle u,v\rangle$, with both vectors normalized. The original comparator supplied the remaining weight. As in the parent, a unique physical match to the emitted verb took precedence. The objectives combined contrastive source--question matching, literal--description agreement and regression to the training action vector.

From 16,632 Kyoto training contexts and 507 action labels, 425 actions had eligible WordNet descriptions or alternatives (7,423 associations).\footnote{\url{https://wordnet.princeton.edu/}} We excluded the held description strings from the additional supervision. In the control we permuted described questions among sampled contexts while retaining the literal questions and shared optical initialization. On one RTX 4090 D the word encoder used 10,000 updates and 19.0 seconds. Each comparator used 2,200 updates, taking 14.2 and 13.2 seconds for aligned and exchanged descriptions respectively. These times exclude parent training, teacher construction, data preparation and evaluation. The complete predictor contains 3,526,604 parameters. Its inference inputs are page and question pixels.

The first wording test used 53 reserved excerpts from Webster's 1913 dictionary. For literary transfer we used published translations to annotate object heads and acceptable English readings before scoring.\footnote{Pu Songling, \textit{Liaozhai zhiyi}, \url{https://www.gutenberg.org/ebooks/51828}. H. A. Giles, trans., \textit{Strange Stories from a Chinese Studio}, \url{https://www.gutenberg.org/ebooks/43629}.} We checked the rendered source page against inherited syntax training and the padded 17-cell verb and object windows against the inherited writer and reader records. The Pu passages were chosen after the Tang overlap audit, with weights already frozen. All 480 Pu responses were generated through the image interface and rescored from saved bitmaps, stopping decisions and source positions. Acceptable answers were registered by the model authorship workflow without independent human adjudication. A single-cell object head represents compounds such as the word for peach. Lean verifies a structural restriction of this decoder: a question-independent pointer for a selected line cannot select two distinct object positions in that line. This restriction matters when a passage contains more than one action.
